% TOP_PROBLEM: {"id": 30007249, "title": "The optimal randomized-versus-quantum query exponent for total Boolean functions", "category_id": 16, "category": {"id": 16, "name": "physics", "display_name": "Mathematical Physics", "description": "Problems at the intersection of mathematics and physics.", "slug": "physics", "order_index": 16, "created_at": "2026-07-31T15:26:25.671Z"}, "status": "open"}
% FIRST_POSTED: 2026-10-06
% ATTEMPT_STATUS: solved
% OPENAI_CLAIM: {"schema_version":1,"source_repo":"https://github.com/openai/math","source_commit":"adc7f1241b42e322a6451854ab7e4b4c146bf78a","release_date":"2026-10-06","match":"full","resolution":"proved","independently_reviewed":false,"second_pass":"agreed","families":[{"family":"284","title":"The optimal quartic separation between randomized and quantum queries","subject":"Mathematical physics","manuscripts":[{"title":"A Nearly Quartic Separation Between Randomized and Quantum Query Complexity","pdf_path":"preprints/A-Nearly-Quartic-Separation-Between-Randomized-and-Quantum-Query-Complexity-October-5-2026/quartic-query-separation.pdf","readme_path":"preprints/A-Nearly-Quartic-Separation-Between-Randomized-and-Quantum-Query-Complexity-October-5-2026/README.md","date":"2026-10-05","theorem_ref":"thm:main (sections/01-introduction.tex:29)","statement_sources":[{"path":"preprints/A-Nearly-Quartic-Separation-Between-Randomized-and-Quantum-Query-Complexity-October-5-2026/build/sections/01-introduction.tex","line":29,"label":"thm:main"}]}],"lean":null,"reasoning_trace":null}]}
\documentclass[11pt]{article}
\usepackage[margin=25mm]{geometry}
\usepackage{fontspec}
\setmainfont{Latin Modern Roman}
\usepackage{amsmath,amssymb,mathtools}
\usepackage{unicode-math}
\setmathfont{Latin Modern Math}
\usepackage{enumitem}
\usepackage{xurl,hyperref,bookmark}
\hypersetup{colorlinks=true,urlcolor=blue}
\setcounter{secnumdepth}{0}
\setlength{\parindent}{0pt}
\setlength{\parskip}{5pt}
\setlength{\emergencystretch}{3em}
\begin{document}
\section*{The optimal randomized-versus-quantum query exponent for total Boolean functions}

\subsection{OpenAI's claim}
OpenAI presents a proof in the selected manuscript. Total Boolean functions are claimed to separate bounded-error randomized and quantum queries by nearly the fourth power, excluding every universal exponent below four.

\subsection{Mathematical statement of the claimed result}
\paragraph{Result 1.} For total Boolean functions, let \(R(f)\) and \(Q(f)\) be bounded-error randomized and quantum query complexities. The source claims the optimal universal power is four:
\[
 R(f)=O((1+Q(f))^4).
\]
For every fixed integer \(k\ge2\), it supplies total functions \(F_{k,m}\), for sufficiently large powers of four \(m\), such that
\[
 Q(F_{k,m})\le C_k\sqrt m(\log m)^{b_k},\qquad
 R(F_{k,m})\ge c_km^{2-1/k}/(\log m)^2.
\]
Letting \(k\) grow excludes any smaller universal exponent.

\subsection{Outline of the argument}
An addressed certificate with outward pointers totalizes a hard partial function. An exact quantum subroutine and certificate checking give the quantum upper bound, while randomized information and pointer constraints force the larger worst-case bit-query cost.

\subsection{Scope relative to this problem}
Total Boolean functions are claimed to separate bounded-error randomized and quantum queries by nearly the fourth power, excluding every universal exponent below four. The question retains these same parameters and conventions. Within these assumptions, this is a claimed affirmative answer to the stated question.

\subsection{Formal verification}
Family 284 has no Lean scope page in this release. Its selected result is not presented as formalized. This site has not executed the Lean code. The scope descriptions do not certify the remaining manuscript arguments.

\subsection{Status caveat}
OpenAI attributes these manuscripts to work by an internal model. Its release cautions that arguments without formalization can contain mistakes. This catalogue has not independently checked the proofs; the labels report the mathematical claims attributed to OpenAI.

\subsection{Sources}
\begin{itemize}
\item OpenAI, \emph{A Nearly Quartic Separation Between Randomized and Quantum Query Complexity}, 2026-10-05; citation key \texttt{OAI:A-Nearly-Quartic-Separation-Between-Randomized-and-Quantum-Query-Complexity-October-5-2026}. \href{https://github.com/openai/math/blob/adc7f1241b42e322a6451854ab7e4b4c146bf78a/preprints/A-Nearly-Quartic-Separation-Between-Randomized-and-Quantum-Query-Complexity-October-5-2026/quartic-query-separation.pdf}{PDF (GitHub)}; \href{https://github.com/openai/math/raw/adc7f1241b42e322a6451854ab7e4b4c146bf78a/preprints/A-Nearly-Quartic-Separation-Between-Randomized-and-Quantum-Query-Complexity-October-5-2026/quartic-query-separation.pdf}{PDF (direct)}; \href{https://github.com/openai/math/blob/adc7f1241b42e322a6451854ab7e4b4c146bf78a/preprints/A-Nearly-Quartic-Separation-Between-Randomized-and-Quantum-Query-Complexity-October-5-2026/README.md}{Manuscript page and citation}; \href{https://github.com/openai/math/blob/main/preprints/A-Nearly-Quartic-Separation-Between-Randomized-and-Quantum-Query-Complexity-October-5-2026/quartic-query-separation.pdf}{Latest version}; Statement source: \href{https://github.com/openai/math/blob/adc7f1241b42e322a6451854ab7e4b4c146bf78a/preprints/A-Nearly-Quartic-Separation-Between-Randomized-and-Quantum-Query-Complexity-October-5-2026/build/sections/01-introduction.tex\#L29}{thm:main}.
\item \href{https://github.com/openai/math/blob/adc7f1241b42e322a6451854ab7e4b4c146bf78a/CONTENTS.md}{OpenAI catalogue entry}.
\end{itemize}
\end{document}
